\documentclass{article}

\usepackage{neurips_2024}

\title{Reinforcement Learning in Pac-Man}

\author{
Anonymous Author(s) \
Affiliation \
Address \
\texttt{email}
}

\begin{document}

\maketitle

\begin{abstract}
This project investigates the effectiveness of various reinforcement learning algorithms in the dynamic environment of Pac-Man. We aim to identify the optimal algorithm by evaluating an agent's adaptability to challenges such as faster ghosts and regenerating pellets.By analyzing performance in this modified version of the typical Pac-Man game, we seek to determine which algorithms are most effective at balancing exploration, adaptability, and long-term reward optimization.
\end{abstract}

\section{Introduction}

% not sure an introduction to the introduction is really necessary 
\iffalse
This report details our COMP4010 project on the application and evaluation of reinforcement learning algorithms within the classic arcade game, Pac-Man.
\fi 

\subsection{Problem}

The central problem this project addresses is identifying the most effective reinforcement learning algorithm for an agent playing Pac-Man under a variety of non-static conditions. The core investigation focuses on the agent's ability to adapt its strategy in response to several environmental challenges. These challenges include:

\begin{itemize}
\item Navigating the maze against a single ghost that moves at a higher speed than the Pac-Man agent.
\item Scenarios where pellets regenerate after being consumed, maintaining a set minimum number of pellets in the maze at any given time.
\end{itemize}


\subsection{Importance}

The game of Pac-Man, while seemingly simple, provides a rich and challenging environment that serves as a classic benchmark for reinforcement learning algorithms. Its clearly defined rules, states, and actions make it an ideal testbed for modeling and agent-based learning. The problem is non-trivial; the optimal strategy is not immediately obvious, requiring a balance between offensive (pellet-seeking) and defensive (ghost-avoiding) tactics.

By studying how different algorithms perform in a dynamic Pac-Man environment, we can gain valuable insights into their robustness, adaptability, and efficiency. These findings have broader implications beyond the game itself, informing the application of reinforcement learning to real-world problems characterized by uncertainty and changing conditions, such as robotics, autonomous navigation, and resource management.

\subsection{MDP Specification}

To formally model the Pac-Man environment for our reinforcement learning agent, we define it as a Markov Decision Process (MDP) with the following specifications:

\begin{itemize}
\item \textbf{States (S):} A state is a complete description of the environment at a given time. It is defined by a $n\times n$ matrix, where $n$ is the width of the board. The board contains information about where walls are, where pellets are, and the locations of Pac-Man and the ghost. 

\item \textbf{Actions (A):} The action space is discrete and consists of the four  directions: {Up, Down, Left, Right}. The agent chooses one of these actions at each time step.
\item \textbf{Rewards (R):} The agent receives a numerical reward based on the outcome of its actions. The reward function is structured as follows:
\begin{itemize}
\item A small positive reward for consuming a standard pellet.
\item A large negative reward (penalty) for being caught by a ghost.
\item A small negative reward for each time step to encourage efficient navigation and discourage stalling.
\end{itemize}
\item \textbf{Transition Model (P):} The transition model P(s',s,a) defines the probability of transitioning from state s to state s' after taking action a. For Pac-Man's movement, the transition is deterministic. The movement of the ghost is also deterministic, follow a policy under which it tries to move towards Pac-Man to minimize the distance between the two. The transition is stochastic in the case of pellet regeneration, which occurs by the random repopulation of empty squares with pellets. 

\item \textbf{Discount Factor ($\gamma$):} A discount factor, $\gamma \in [0,1]$, is used to weigh the importance of future rewards. A value closer to 1 results in an agent that is more far-sighted, while a value closer to 0 results in an agent that prioritizes immediate rewards. In the case of Pac-Man, where we look to get a high score in the game, $\gamma$ should be relatively high to ensure the agent puts emphasis on future rewards. 
\end{itemize}

\section{Approaches}

\subsection{Previous Approaches}

\subsection{Algorithms Used}

In this project, we implement the algorithms SARSA and Q-Learning. 

\section{Empirical Studies}

\subsection{Experiment Procedure}

\subsection{Results}

\subsubsection{Q-Learning Performance and Behaviour}

We trained the Q-learning agent for 3000 episodes using optimistic initialization, an $\epsilon$-greedy policy with exponential decay, and a constant learning rate. Performance was evaluated every 100 episodes using near-greedy roll-outs ($\epsilon \approx 0$). The results show a clear progression of learning followed by late-stage instability once exploration diminished.

\paragraph{Early Training (Episodes 1--1000).}
During the initial phase, Pac-Man survived only 10--20 steps per episode, collecting approximately 10--15 pellets. Returns were low (4--9), and evaluation roll-outs showed no pellets collected. Although pellets are plentiful on the board, the agent's early behaviour was dominated by random exploration and frequent deaths, which limited the amount of useful reward collected per episode and slowed the learning process.

\paragraph{Mid Training (Episodes 1000--2000).}
The agent began to demonstrate structured behaviour, consistently collecting 14--20 pellets and surviving 15--25 steps. Pac-Man reacted to ghost proximity, avoided immediate capture, and reduced purely random movement. This indicated that the Q-values were beginning to encode short-term survival strategies.

\paragraph{Breakthrough phase (Episodes 2000-2500).}
A significant improvement occurred around episode 2000. Pac-Man collected 50--75 pellets, survived 60--120 steps, and achieved returns between 40 and 65. Evaluation metrics showed similar gains. The agent dropped a cautious strategy focused on maintaining distance from the ghost while clearing large clusters of pellets. This reflects successful propagation of long-term value estimates in the Q-table.

\paragraph{Late Training instability (Episodes 2500+).}
After $\epsilon$ decayed to its minimum of 0.08, performance became more volatile. With limited exploration and a constant learning rate, the agent over-fitted to deterministic behavioral loops and occasionally reverted to less effective patterns. Pellet regeneration also introduced additional variability by altering the environment mid-episode.

\paragraph{Qualitative Behavioral Observations.}
During evaluation, we observed that Pac-Man learned to exploit weaknesses in the ghost's Manhattan-distance chase policy. In particular, the agent frequently positioned itself behind walls where the ghost repeatedly moved in circles attempting to minimize the direct grid distance, failing to path around obstacles. This behaviour does not appear directly in training metrics but demonstrates that Q-learning successfully identified structural asymmetries in the environment and used them to extend survival time.

\paragraph{Summary.}
Overall, Q-learning demonstrated meaningful learning: returns increased from approximately 5 to over 60, and pellet collection improved from ~12 to ~75. While the final policy was not globally optimal and displayed sensitivity to hyper-parameters and deterministic transitions, it consistently outperformed random behaviour and exhibited non-trivial strategic decision-making.

\section{Conclusion}

\subsection{Future Work}

\end{document}